\documentclass{article}
\usepackage{amsmath}
\begin{document}
\section{Oct 20-22}
1,2,4,6,7

1:$\Delta \nu = \frac{c}{2\mu L}$ $ = \frac{3*10^8}{1}=3*10^8$ Bec : $\Delta \nu = 6*10^8 $ 

So : $n=\frac{\Delta \nu}{3*10^8}$=2.

q could be : q-1;q;q+1.

and
$\nu = \frac{qc}{2\mu L}$ so $ q=\frac{2\mu L\nu}{c}$ $=\frac{5.85*10^{14}}{3*10^8}=1.95*10^{16}$and q+1 q-1.

2:

$\Delta \nu = \frac{3*10^8}{2}=1.5*10^8$=$0.1*\Delta \nu$

so n=10 ,q could have 11 . if q only neede 1 ,so 

$\Delta \nu = 1.5*10^9=\frac{3*10^8}{2L}$so L=0.1 m .

4:

(1):$P/S=\frac{50W}{(25\mu m)^2\pi}{\mu m}^2$= K

(2):$\frac{K}{10^4}$=$N_1$;$\frac{K}{10^3}$=$N_2$

6:

$\theta=\frac{1.22*632.8nm}{2mm}=\theta_1$

$\theta_0$=$\frac{\lambda}{\pi \omega_0}$=$\frac{\lambda}{\pi \sqrt\frac{\lambda L}{2\pi}}$

7:

$\omega_0$=$\sqrt\frac{\lambda L}{2\pi}$

$\omega (z)$=$\omega_0\sqrt{1+{\frac{\lambda z}{\pi \omega_0 ^2}^2}}$

$L=0.4;\lambda=0.6328\mu m ;z=56 cm $

9:

$R=z[1+(\frac fz)^2]$$\frac{3\sqrt3}7=f,R_1'=R_2'=\frac{6\sqrt3}{7}$

$\omega_0=\sqrt\frac{\lambda f}{\pi}=\sqrt\frac{\lambda 3\sqrt3}{7\pi}$

$\omega=\omega_0\sqrt{1+(\frac fz)^2}=\sqrt\frac{\lambda 3\sqrt3}{7\pi}\sqrt{1+(\frac {3\sqrt3}{7z} )^2}$

$\frac{R(z_1)-z_1}{z_1}=(\frac f{z_1})^2;\frac{R(z_2)-z_2}{z_2}=(\frac f{z_2})^2$

$\frac{1.5-z_1}{3+z_1}=\frac{1-z_1}{z_1};z_1+z_2=1;\frac67=z_1;z_2=\frac17;$

10:

$2=1[1+ f^2]$;f=1;$\omega=\omega_0\sqrt{1+(\frac fz)^2}=\sqrt{\frac{2\lambda}{\pi}}=2.5977mm$

$\omega=\omega_0\sqrt{1+(\frac fz)^2}=\omega_0\sqrt{1+f^2}=\omega_0\sqrt{R}=\sqrt{\frac{\lambda\sqrt{R-1}R}{\pi}}=0.3$

$0.3^2\pi/\lambda=R\sqrt{R-1}$;R=41.7733cm.

$\omega_0=\sqrt{\frac{\lambda \sqrt R}\pi}$

$\omega=\omega_0\sqrt{1+f^2};R=1+f^2;f=\sqrt{R-1};\omega_0\sqrt{R}=\sqrt{\frac{\lambda R\sqrt R}\pi} $

$  R=(0.3^2\pi/\lambda)^{2/3}=$

$f=\frac{\sqrt{2(R-2)^2(2R-2)}}{2R-4}=\frac{2\sqrt{R-1}}{2}=\sqrt{R-1}$

$\omega_0=\sqrt{\frac{\lambda \sqrt{Rz}}{\pi}\frac Rz};f=\sqrt{(\omega_0^2\pi/\lambda)^{2/3}}$

$(\omega_0^2\pi/\lambda)^{2/3}=R$

$\omega=\omega_0\sqrt{1+(\omega_0^2\pi/\lambda)^{2/3}}$
\newpage
\section{redo HW all}
redo

\subsection{chapter 1}

2:

(1):

\begin{align*}
  &\frac{n_2}{n_1}=\frac{g_2}{g_1}e^{\frac{-h\nu}{kT}}\\
  &=e^{-\frac{h\times3\times10^{9}}{1.38\times10^{-23}\times300}}=0.995\\
\end{align*}

(2)

\begin{align*}
  &0.1=e^{-\frac{h\frac c{\lambda}}{1.38\times10^{-23}T}}\\
  &-ln10=-\frac{h\frac c{\lambda}}{1.38\times10^{-23}T}\\
  &T=\frac{h\frac c{\lambda}}{1.38\times10^{-23}ln10}=6259 k\\
\end{align*}
3:

(1):

\begin{align*}
&\frac{n_2}{n_1}=\frac{g_2}{g_1}e^{-\frac{E}{kT}}\\
&n_2=n_1\frac{g_2}{g_1}e^{-\frac{E}{kT}}=10^{20}\times4\times e^{-\frac{1.64\times10^{-18}}{1.38\times10^{-34}\times2700}}=30.66\\
\end{align*}

(2):

\begin{align*}
  P=I \times A=h\nu \times n=1.64 \times 10^{-18} \times 10^8 \times n_2=4.92 \times 10^{-9}\\
\end{align*}

4:

(1):

\begin{align*}
  &\frac{q}{q_n}=\frac{B_{21}\rho}{A_{21}}=\frac{c^3\rho}{8\pi h\nu^3}=\frac1{2000}\\
  &\rho=\frac{8\pi h\nu^3}{2000 \times c^3}=\frac{8\pi h}{2000 \times \lambda^3}=3.85 \times 10^{-17} \\
\end{align*}

(2):

\begin{align*}
  &\frac{q}{q_s}=\frac{c^3\rho}{8\pi h\nu^3}=\frac{\lambda^3\rho}{8\pi h}=7.603 \times 10^9\\
\end{align*}

11:

\begin{align*}
&\nu=\nu_0(1+\frac vc);\nu_{0.1}=\nu_0 \times 1.1;\nu_{-0.1}=\nu_0 \times 0.9\\
&\nu_{0.5}=\nu_0 \times 1.5;\nu_{-0.5}=\nu_0 \times 0.5\\
&\nu_0=4.74 \times 10^{14}\\
&\nu_{0.1}=5.214 \times 10^{14};\nu_{-0.1}=4.266 \times 10^{14};\\
&\nu_{0.5}=7.11 \times 10^{14};\nu_{-0.5}=2.370 \times 10^{14}\\
\end{align*}

12:

\begin{align*}
 & \delta \nu=\nu-\nu_0=\nu_0 \times (1+\frac vc)-\nu_0= \frac vc\times \nu_0=8.848 \times 10^{8}Hz\\
\end{align*}

13:

(1):

\begin{align*}
&I=I_0\times e^{-Gz};\frac{I}{I_0}=e^{-Gz}=e^{-0.01 \times 10^3 \times 10 \times 10^{-2}}=e^{-1}=36.78\%\\
\end{align*}

(2):

\begin{align*}
&2I=Ie^{G \times 1m};ln2=G\times 1;G=0.6931\\
\end{align*}

\subsection{chapter 2}

1:

\begin{align*}
  &G= \Delta h\nu/cf(\nu)B_{21}=5 \times 10^{24} \times  f(\nu)\frac{A_{21}c^2}{8\pi \nu^2\mu^2}=5 \times 10^{24}\frac1{2 \times 10^{11}}\frac{\lambda^2}{8\pi\tau\mu^2}\\
  &=71.037\\
\end{align*}

2:

\begin{align*}
&n_1+n_2=10^{18};n_2=4n_1;n_1=2 \times 10^{17};n_2=8 \times 10^{17};\Delta n=n_2-n_1\frac{g_2}{g_1}=\\
&(8 \times -2 \times \frac{3}5) \times 10^{17}=\frac{34}5 \times 10^{17}\\
&g= \Delta n\frac{h\nu}{c}f(\nu)B_{21}=\frac{34}5 \times 10^{17} \frac{h\nu}cf(\nu)A_{21}\frac{c^3}{8\pi h\nu^3}=\frac{34}5 \times 10^{17}f(\nu)\frac{c^2}{8\pi\nu^2\tau}\\
&=72.228\\
\end{align*}

3:

(a)

\begin{align*}
&0<(1-\frac{L}{R})^2<1\\
&-1<1-\frac LR<1:\frac LR>0;\frac LR<2,R>0,R>\frac L2\\
&R>30cm\\
\end{align*}

(b):

\begin{align*}
&0<(1-\frac{L}{4L})(1-\frac LR)<1;0<\frac34(1-\frac LR)<1\\
&\frac LR<1;-\frac13<\frac LR \\
&R>0:R>L;R>-3L\\
&R<0:R<L:R<-3L\\
&R:R>L | R<-3L\\
\end{align*}

7:

\begin{align*}
&G=\frac{(\frac{\Delta\nu}{2})^2G_0(\nu_0)}{(\nu-\nu_0)^2+(1+\frac{I}{I_s})(\frac{\Delta\nu}2)^2}\\
&\text{if:}(\nu-\nu_0)^2>>(1+\frac{I}{I_s})(\frac{\Delta\nu}2)^2\\
&|\nu-\nu_0|=\Delta\nu'=\sqrt{1+\frac{I}{I_s}}\Delta\nu\\
&I=I_s;\Delta\nu'=\sqrt{2}\Delta\nu\\
\end{align*}

8:

\begin{align*}
&\sigma_e(\nu)=\frac{G(\nu)}{\Delta n}=\frac{\Delta n\frac{h\nu}cf(\nu)B_{21}}{\Delta n}=\frac{h\nu f(\nu)B_{21}}{c}\\
&=\frac{ f(\nu)A_{21}c^2}{ 8\pi \nu^2}=\frac{f(\nu)c^2}{8\pi\mu^2\tau\nu^2}\\
\end{align*}

10:

\begin{align*}
  (a)&G_D(\nu)=\frac{G_D(\nu_1)}{\sqrt{1+\frac{I}{I_s}}}=\frac{3 \times 10^{-4}}{10^{-3} \times \sqrt{1+\frac{5}{3}}}=0.1837\\
  (b)&a_a=a_i-\frac1{2L}ln(r_1r_2)=G\\
     &a_i-G=\frac1{2L}ln(r_1r_2)=\frac1Lln(r)=(a_i-G)\\
     &=r=e^{L(a_i-G)}=0.9906\\
  (c)&P=tIA=0.008 \times 50 \times 10^{4} \times 0.11 \times 10^{-6}=4.4 \times 10^{-4}\\
     &=44mW\\
\end{align*}

11:

\begin{align*}
&G=\Delta n h\frac{\nu}{c}f(\nu)B_{21}\\
&\Delta n=\frac{cG}{h\nu f(\nu)B_{21}}=\frac{8\pi \nu^2a_a\tau}{ f(\nu)c^2}\\
&R=1-0.1 \times a_a=0.9833;9.833=10- a_a\\
&a_a=0.167\\
&\Delta n=1.048 \times 10^{15}\\
\end{align*}

\subsection{chapter 3}

1:

\begin{align*}
&\Delta \phi=\frac{c}{2\mu L}=3 \times  10^8\\
&\frac{\Delta\nu}{\Delta\phi}=2;N_q=3\\
&\phi=5.85 \times 10^{14}=\frac{qc}{2\mu L}=q \times 3 \times 10^8\\
&q=1.95 \times 10^6;q+1=1.95 \times 10^6+1;q-1=1.95 \times 10^6-1\\
\end{align*}

2:

\begin{align*}
&\Delta\phi=\frac{c}{2\mu L}=1.5 \times 10^8\\
&\frac{\Delta\nu}{\Delta\phi}=10;N_q=11\\
&\text{if:}N_q=1;\Delta\phi>\Delta\nu:\frac{c}{2\mu L}=\frac{1.5 \times 10^8}{L}>1.5 \times 10^9\\
&L<0.1m\\
\end{align*}

4:

(1):

\begin{align*}
&p/s=\frac{50}{s};s=\frac{\pi d^2}4=1.963 \times 10^{-9}\\
&p/s=2.54 \times 10^{10}W/m^2=2.54 \times 10^6 W/cm^2\\
\end{align*}

(2):

\begin{align*}
&254 \times \\
&2540 \times \\
\end{align*}
\newpage
6:

\begin{align*}
&\theta_0=\frac{\lambda}{\pi\omega_0}=\frac{\lambda}\pi\sqrt{\frac{\pi}{\lambda f}}=\sqrt{\frac{\lambda}{\pi f}}=1.1588 \times 10^{-3}\\
&\theta_r=\frac{1.22 \times 632 \times 10^{-9}}{2 \times 10^{-3}}=3.855 \times 10^{-4}\\
\end{align*}

7:

\begin{align*}
&\omega=\omega_0\sqrt{1+(\frac zf)^2};\omega_0=\sqrt{\frac{\lambda f}\pi}=2 \times 10^{-4}\\
&\omega=2 \times 10^{-4} \times \sqrt{1+(\frac{0.56}{0.2})^2}=5.94 \times 10^{-4}\\
\end{align*}

9:

\begin{align*}
&R_1z_1=z_1^2+f^2;R_2z_2=z_2^2+f^2\\
&R_1z_1-z_1^2=R_2z_2-z_2^2;z_1+z_2=1\\
&R_1z_1-z_1^2=R_2(1-z_1)-(1-z_1)^2;R_1=1.5,R_2=4\\
&1.5z_1-z_1^2=4-4z_1-1-z_1^2+2z_1;3.5z_1=3;z_1=\frac67;z_2=\frac17\\
&f^2=\frac47-\frac1{49}=\frac{27}{49};f=\frac{3\sqrt{3}}7\\
&\omega_0=\sqrt{\frac{\lambda f}\pi}=3.486 \times 10^{-4}\\
&\omega_2=\omega_0\sqrt{1+(\frac {z_2}f)^2}=3.549 \times 10^{-4}\\
&\omega_1=\omega_0\sqrt{1+(\frac{z_1}f)^2}=5.32 \times 10^{-4}\\
\end{align*}
\newpage
10:

\begin{align*}
  (a)&\omega_0=\sqrt{\frac{\lambda f}\pi}=1.836 \times 10^{-3}\\
     &\omega=\omega_0\sqrt{1+(\frac{z}f)^2}=2.597 \times 10^{-3}\\
  (b)&\omega=\omega_0\sqrt{1+(\frac zf)^2}=\sqrt{\frac{f\lambda}\pi}\sqrt{1+(\frac zf)^2}=0.3 \times 10^{-2}\\
     &\text{solve f :}f=2.2161\\
     &R=z+\frac{f}{z}^2=1+f^2=5.911\\
\end{align*}
\subsection{chapter 4}

1:

\begin{align*}
&\Delta\phi=\frac{c}{2\mu (L_2+L_3)};N_q=3,\frac{\Delta\nu}{\Delta\phi}=2\\
&750 \times 10^{6}=\frac{3 \times 10^8}{2\mu(L_2+L_3)};(L_2+L_3)=\frac15=0.2\\
\end{align*}

2:

\begin{align*}
  (1)&\omega_{00}=\sqrt{\frac{\lambda f}\pi}=1.419 \times 10^{-4}\text{add a hole at mirror surface}\\
     &\omega=\omega_0\sqrt{1+(\frac zf)^2}=2.007 \times 10^{-4}\\
  (2)&\omega_{11}=\sqrt{2+1}\omega_{00}=3.476 \times 10^{-4}\\
\end{align*}

3:

\begin{align*}
&\omega_0=\sqrt{\frac{f\lambda}\pi};f=\omega_0^2\pi/\lambda=0.1985\\
&q=0.6+0.1985i;T=\begin{bmatrix}1&0\\-\frac1f&1\end{bmatrix}\\
&q'=\frac{q}{-\frac{100}3 \times(0.6+0.1985i)+1}=-0.0314+4.903 \times 10^{-4}i\\
\end{align*}
\begin{align*}
&\omega'=\sqrt{\frac{\lambda f'}\pi}=9.9 \times 10^{-6}\\
\end{align*}

4:

\begin{align*}
&\theta_1=\frac{\lambda}{\pi\omega_1}=1.007 \times 10^{-3}\\
&\theta_2=\frac{\lambda}{\pi\omega_2}=4.026 \times 10^{-3}\\
&\text{4 times bigger}
\end{align*}

5:

\begin{align*}
&M=\frac{f_2\omega}{f_1\omega_0}=\sqrt{2} \times \frac{f_2}{ f_1}=11.31\\
\end{align*}

8:

\begin{align*}
  (1)&\Delta\phi=\frac{c}{2\mu L}=\frac{c}{2 \times 1.5}=10^8\\
     &T=10^{-8}s\\
     &\tau=\frac{T}{N_q}=\frac{10^{-8}}{1000}=10^{-11}\\
     &P_{peak}=NP_{aver}=1000 \times 1=10^3\\
  (2)&T=\frac{2\pi}{\omega}=\frac{2\pi}{2\pi f}=\frac1f=10^{-8}\\
     &f=10^8\\
\end{align*}

9:

\begin{align*}
&N_q=\frac{\Delta\nu}{\Delta\phi}+1=\frac{7.5 \times 10^{12}}{\frac{c}{2(1.52 \times 0.2+0.1)}}+1=20201\\
&P_{peak}=N_q \times P_{aver};\\
&\text{20201 times bigger}
\end{align*}

\end{document}

